\documentclass[11pt,a4paper]{article}

\usepackage[T1]{fontenc}
\usepackage[utf8]{inputenc}
\usepackage[scaled=0.95]{helvet}
\renewcommand{\familydefault}{\sfdefault}
\usepackage{microtype}
\usepackage[a4paper,margin=19mm,headheight=22pt,footskip=12mm]{geometry}
\usepackage[dvipsnames,table]{xcolor}
\usepackage{hyperref}
\usepackage{bookmark}
\usepackage{enumitem}
\usepackage{titlesec}
\usepackage{fancyhdr}
\usepackage[most]{tcolorbox}
\usepackage{parskip}
\usepackage{lastpage}
\usepackage{xurl}

\definecolor{DARESNavy}{HTML}{102A43}
\definecolor{DARESTeal}{HTML}{0B7285}
\definecolor{DARESMuted}{HTML}{66788A}
\definecolor{DARESLine}{HTML}{D9E2EC}
\definecolor{DARESBlueLight}{HTML}{EAF2FF}
\definecolor{DARESTealLight}{HTML}{E7F5F7}
\definecolor{DARESOrangeLight}{HTML}{FFF4E6}
\definecolor{DARESGreenLight}{HTML}{EAF7F0}
\definecolor{DARESPaper}{HTML}{F7F9FC}
\definecolor{DARESCoral}{HTML}{E76F51}

\hypersetup{
  colorlinks=true,
  linkcolor=DARESTeal,
  urlcolor=DARESTeal,
  citecolor=DARESTeal,
  pdftitle={Unmanned Aircraft Control System - Learning and Project Documentation Guide},
  pdfauthor={DARES Electronics},
  pdfsubject={Fixed-wing guidance, navigation and control onboarding, simulation, and documentation},
  pdfkeywords={UAV, fixed-wing, flight control, GNC, PX4, ArduPilot, SITL, HIL, state estimation, TECS, L1}
}
\urlstyle{same}
\setcounter{tocdepth}{1}
\setlength{\parindent}{0pt}
\setlength{\parskip}{5pt}
\setlist[itemize]{leftmargin=5.5mm,itemsep=2.5pt,topsep=3pt}
\setlist[enumerate]{leftmargin=7mm,itemsep=3pt,topsep=3pt}
\emergencystretch=3em

\titleformat{\section}
  {\Large\bfseries\color{DARESNavy}}
  {\thesection.}{0.55em}{}
  [\vspace{2pt}\color{DARESTeal}\titlerule]
\titleformat{\subsection}
  {\large\bfseries\color{DARESTeal}}
  {\thesubsection}{0.5em}{}
\titleformat{\subsubsection}
  {\normalsize\bfseries\color{DARESNavy}}
  {\thesubsubsection}{0.5em}{}
\titlespacing*{\section}{0pt}{14pt}{8pt}
\titlespacing*{\subsection}{0pt}{10pt}{3pt}
\titlespacing*{\subsubsection}{0pt}{7pt}{2pt}

\pagestyle{fancy}
\fancyhf{}
\lhead{\small\bfseries\color{DARESTeal}DARES ELECTRONICS}
\rhead{\small\color{DARESMuted}Unmanned Aircraft Control System}
\lfoot{\scriptsize\color{DARESMuted}Learning and project documentation guide}
\rfoot{\scriptsize\color{DARESMuted}Page \thepage\ of \pageref*{LastPage}}
\renewcommand{\headrulewidth}{0.4pt}
\renewcommand{\footrulewidth}{0.4pt}
\renewcommand{\headrule}{\hbox to\headwidth{\color{DARESLine}\leaders\hrule height \headrulewidth\hfill}}
\renewcommand{\footrule}{\hbox to\headwidth{\color{DARESLine}\leaders\hrule height \footrulewidth\hfill}}

\newtcolorbox{callout}[1]{
  enhanced,
  breakable,
  colback=DARESTealLight,
  colframe=DARESLine,
  boxrule=0.4pt,
  borderline west={3pt}{0pt}{DARESTeal},
  arc=0pt,
  left=3mm,
  right=3mm,
  top=2.5mm,
  bottom=2.5mm,
  title={#1},
  coltitle=DARESNavy,
  fonttitle=\bfseries,
  attach title to upper={\par\vspace{1mm}}
}

\newtcolorbox{warningbox}[1]{
  enhanced,
  breakable,
  colback=DARESOrangeLight,
  colframe=DARESLine,
  boxrule=0.4pt,
  borderline west={3pt}{0pt}{DARESCoral},
  arc=0pt,
  left=3mm,
  right=3mm,
  top=2.5mm,
  bottom=2.5mm,
  title={#1},
  coltitle=DARESNavy,
  fonttitle=\bfseries,
  attach title to upper={\par\vspace{1mm}}
}

\newtcolorbox{resourcebox}{
  enhanced,
  breakable,
  colback=DARESBlueLight,
  colframe=DARESLine,
  boxrule=0.4pt,
  arc=0pt,
  left=3mm,
  right=3mm,
  top=2.5mm,
  bottom=2.5mm
}

\newcommand{\Resource}[3]{%
  \begin{resourcebox}
    \href{#2}{\textbf{\color{DARESTeal}#1}}\par
    {\small #3}
  \end{resourcebox}
}

\newcommand{\RepoEntry}[2]{%
  \texttt{\detokenize{#1}}\hfill{\small\color{DARESMuted}#2}\par
}

\begin{document}
\hypersetup{pageanchor=false}

\begin{titlepage}
  \thispagestyle{empty}
  \pagecolor{DARESPaper}
  \color{DARESNavy}

  \vspace*{20mm}
  {\large\bfseries\color{DARESTeal}DARES ELECTRONICS\par}
  \vspace{12mm}

  {\fontsize{30}{34}\selectfont\bfseries Unmanned Aircraft\par}
  \vspace{1mm}
  {\fontsize{30}{34}\selectfont\bfseries Control System\par}
  \vspace{3mm}
  {\Large\color{DARESMuted}Learning and Project Documentation Guide\par}

  \vspace{15mm}
  \begin{tcolorbox}[
    colback=DARESNavy,
    colframe=DARESNavy,
    arc=0pt,
    boxrule=0pt,
    left=4mm,right=4mm,top=4mm,bottom=4mm]
    \centering\bfseries\color{white}
    PRINCIPLE \(\rightarrow\) MODEL \(\rightarrow\) SIMULATION
    \(\rightarrow\) REFERENCE STACK \(\rightarrow\) HIL \(\rightarrow\) FLIGHT TEST
    \(\rightarrow\) MISSION
  \end{tcolorbox}

  \vspace{15mm}
  {\large\color{DARESMuted}
    A beginner-friendly reference for learning fixed-wing flight mechanics,
    building and validating a six-degree-of-freedom model, bringing up a
    reference autopilot stack, documenting engineering work, and progressing
    from stabilised flight toward guidance, mission autonomy and in-house
    control functions.
    \par}

  \vfill
  \begin{callout}{CURRENT PRIORITY}
    Fixed-wing flight mechanics \(\rightarrow\) six-degree-of-freedom model
    \(\rightarrow\) trim and linearisation \(\rightarrow\) SITL
    \(\rightarrow\) rate and attitude loops \(\rightarrow\) speed and height
    control \(\rightarrow\) path following \(\rightarrow\) HIL
    \(\rightarrow\) first flight. Use PX4 or ArduPilot as the reference stack to
    connect control theory to a flying aircraft before developing in-house
    guidance and control functions.
  \end{callout}

  \vfill
  {\color{DARESCoral}\rule{32mm}{2pt}\par}
  \vspace{3mm}
  {\small\color{DARESMuted}Version 1.0 \textbar\ September 2026\par}
\end{titlepage}
\nopagecolor
\hypersetup{pageanchor=true}

\tableofcontents
\clearpage

\section{Purpose and scope}

This guide gives new Unmanned Aircraft Control System (UACS) members a shared
starting point. It explains what to learn, what to model and simulate, what to
configure on a reference autopilot stack, what to document, and what evidence
is required before the team treats a control configuration as validated on the
team's competition fixed-wing platform.

\subsection{Learning objectives}

By using this guide, a member should be able to:

\begin{itemize}
  \item Explain how a fixed-wing aircraft produces lift and drag, how its moments arise, and why it needs continuous control rather than only a throttle command.
  \item Describe the reference frames, the state variables and the attitude representations a flight controller actually uses.
  \item Build or review a six-degree-of-freedom model, trim it, linearise it and interpret its longitudinal and lateral modes.
  \item Explain how sensors and a state estimator produce the estimate that every control loop depends on.
  \item Describe the cascade from guidance to speed and height, attitude and angular-rate control, and how control allocation maps commands onto actuators.
  \item Bring up a reference autopilot stack in software-in-the-loop simulation, then hardware-in-the-loop, then on a bench, then in a first flight.
  \item Document models, simulation results, parameter sets, bench tests and flight tests so that another member can reproduce them.
  \item Recognise when the aircraft is ready for mission modes, higher-level autonomy and in-house control functions.
\end{itemize}

\begin{callout}{IMPORTANT BOUNDARY}
  This guide is an onboarding and documentation standard. It does not replace
  the competition rules, airframe and component datasheets, the team's safety
  and flight-approval procedures, airworthiness and radio requirements, or
  validation on the actual aircraft. Where this guide and an approved procedure
  disagree, the approved procedure wins.
\end{callout}

\subsection{Scope and interfaces}

This guide covers the aircraft-level chain: state estimation, guidance,
control, control allocation, flight modes and failsafe behaviour. It does not
cover the design of the motor, the ESC, the battery-management system or
vision perception. Those are separate workstreams with their own documentation,
and the interfaces to them are listed in Section 9.

\begin{callout}{SCOPE RULE}
  A control-system member is expected to understand the interface to the other
  subteams, not to redesign their hardware. If a control requirement implies a
  change in the propulsion, power or perception interface, raise it as an
  interface change with the responsible subteam rather than working around it
  inside the flight controller.
\end{callout}

\subsection{Recommended working method}

Learn enough flight mechanics to understand the physical system, build a model
that can be explained block by block, validate the control structure in
simulation, bring up the reference stack on the bench, compare predicted and
measured behaviour, and only then develop in-house functions. Modelling,
simulation, stack configuration, bench testing and flight testing are one
chain: each step should create the evidence the next step depends on. Do not
treat them as separate projects, and do not replace a missing measurement with
a better-looking plot.

\section{Recommended learning path}

The sequence below is designed for members with no previous flight-control
experience. Complete each stage well enough to explain it in your own words
before moving to the next stage.

\subsection{Why a fixed-wing aircraft needs a control system}

A fixed-wing aircraft generates lift from airflow over a wing, and it flies
only while airspeed is above the stall speed. Its motion is described by forces
and moments in three axes. A wing alone is usually not stable enough, or not
stable in the way the mission requires, so the flight controller continuously
adjusts the control surfaces and thrust to hold a commanded attitude, speed,
height and path. Unlike a multirotor, a fixed-wing aircraft cannot hover, must
be launched and recovered, and couples speed and height strongly: pulling up
trades speed for height, and adding power changes both.

\textbf{Learning check}
\begin{itemize}
  \item Explain lift, drag, weight and thrust in straight and level flight, and what changes in a climb or a bank.
  \item Explain why stall speed matters for every control decision, including failsafe behaviour.
  \item Explain the roles of aileron, elevator, rudder and throttle.
  \item Explain why a fixed-wing aircraft cannot be stabilised by thrust alone.
\end{itemize}

\subsection{Reference frames, state and attitude}

A flight controller works with several frames that must not be mixed up:
earth-fixed north-east-down, body-fixed forward-right-down, and the wind and
stability frames used for aerodynamic angles. The state normally includes
position, velocity, attitude and angular rate, plus the derived angles of
attack and sideslip. Attitude is commonly carried as a quaternion inside the
estimator and controller, and shown to a human as Euler angles.

\textbf{Learning check}
\begin{itemize}
  \item Write down what each axis of the body frame means and which sign convention is used.
  \item Convert between Euler angles and a rotation matrix or quaternion, and explain gimbal lock.
  \item Explain the difference between ground speed, airspeed and wind.
  \item Explain why an estimator cannot simply integrate gyroscope data for more than a few seconds.
\end{itemize}

\subsection{Airframe model, trim and linearisation}

The airframe model relates the state and the control inputs to forces and
moments, usually through non-dimensional aerodynamic coefficients. The model
must be trimmed at an operating point before it means anything: trim gives the
steady control settings for a chosen speed and height. Linearising around trim
separates the motion into a longitudinal set and a lateral-directional set, and
exposes the modes a controller has to handle, such as the short period, the
phugoid, the roll subsidence, the spiral and the Dutch roll.

\textbf{Learning check}
\begin{itemize}
  \item Explain what each aerodynamic coefficient contributes to the motion, and which ones dominate stability.
  \item Produce a trim solution and explain why it is not unique.
  \item Identify the aircraft's modes from a linear model and say which ones are poorly damped.
  \item Explain what a change in centre of gravity, mass or airspeed does to trim and to stability.
\end{itemize}

\subsection{Sensors and state estimation}

The estimator turns noisy, biased and differently-rate sensor data into the
state the controller uses. Typical inputs are an inertial measurement unit, a
magnetometer, a barometer, a GNSS receiver and, on a fixed-wing aircraft, a
pitot-static airspeed sensor. A complementary or Mahony-style filter is the
right first step for attitude because it can be explained and tested; a full
extended Kalman filter adds multi-sensor fusion, bias estimation and health
reporting, at the cost of complexity.

\textbf{Learning check}
\begin{itemize}
  \item Explain what each sensor measures, its update rate, its noise, and its failure mode.
  \item Explain how accelerometer and gyroscope data are combined to estimate attitude.
  \item Explain why airspeed matters more on a fixed-wing aircraft than on a multirotor.
  \item Explain what estimator health means and why the controller must react when it degrades.
\end{itemize}

\Resource
  {Kalman and Bayesian Filters in Python (free online book)}
  {https://github.com/rlabbe/Kalman-and-Bayesian-Filters-in-Python}
  {Practical, code-first introduction to filters from complementary filtering through to extended and unscented Kalman filters.}

\Resource
  {PX4: ECL EKF tuning and estimator health}
  {https://docs.px4.io/main/en/advanced_config/tuning_the_ecl_ekf.html}
  {How a production estimator is configured, how innovation and health checks are interpreted, and what commonly goes wrong.}

\subsection{Control loops: from angular rate to trajectory}

Flight control is a cascade. The fastest loop controls angular rate; the next
loop controls attitude; above that, speed and height are controlled together
because they are coupled; the outer loop follows a path or a mission. A rate
loop is damped and well conditioned, an attitude loop gives the pilot or
guidance layer a predictable aircraft, and the speed and height layer manages
energy. Total energy control is a common way to formalise that trade. Path
following closes the loop onto a line or an arc.

\textbf{Learning check}
\begin{itemize}
  \item Explain why the loops are nested and what sets the bandwidth of each one.
  \item Explain proportional, integral and derivative action, and why an integral term needs anti-windup limits.
  \item Explain what total energy control does and why throttle and elevator must be coordinated.
  \item Explain how a path-following law uses a look-ahead distance and what happens when it is too short or too long.
\end{itemize}

\Resource
  {PX4: controller diagrams}
  {https://docs.px4.io/main/en/flight_stack/controller_diagrams.html}
  {The full fixed-wing control cascade as implemented in a production open-source autopilot, from rate loop to total energy and path following.}

\Resource
  {ArduPilot: total energy control system tuning guide}
  {https://ardupilot.org/plane/docs/tecs-total-energy-control-system-for-speed-height-tuning-guide.html}
  {Explains the speed and height control problem, the parameters that shape it, and the symptoms of a badly tuned energy controller.}

\Resource
  {ArduPilot: navigation and path-following tuning}
  {https://ardupilot.org/plane/docs/navigation-tuning.html}
  {Practical treatment of the L1-style path-following controller, tracking quality and the effect of its parameters.}

\subsection{Control allocation and actuators}

Control allocation maps the commands the controller produces onto the physical
actuators: aileron, elevator, rudder, throttle and any flaps or additional
surfaces. It has to respect actuator limits, rates, trim offsets and direction
conventions, and it has to fail in a defined way when an actuator saturates or
is lost. Servo direction, travel, neutral and reversal must be verified on the
bench before any powered test.

\textbf{Learning check}
\begin{itemize}
  \item Explain how a roll command becomes aileron deflection, and what the rudder is doing at the same time.
  \item Explain saturation, rate limiting and how a controller should behave when an actuator reaches its limit.
  \item Explain how a misconfigured servo direction or reversed elevator would appear in flight, and how bench testing catches it.
  \item Explain what changes if the aircraft is a hybrid that also has vertical-lift motors.
\end{itemize}

\section{Simulation before flight}

Simulation is used to test understanding, tune a control structure and reduce
risk before any hardware is energised. It is not a substitute for measurement.
The first model should be simple enough that every block and every parameter
can be explained by the member who built it.

\subsection{Minimum first simulation}

\begin{itemize}
  \item A six-degree-of-freedom airframe model with documented mass, inertia, geometry and aerodynamic coefficients.
  \item A stated source for every aerodynamic coefficient, with an explicit note where a value is estimated or borrowed from a similar aircraft.
  \item Actuator models with documented limits, rates and time constants.
  \item Sensor models with realistic noise, bias, latency and update rates.
  \item Atmosphere, wind and turbulence inputs, including a steady crosswind case.
  \item A trim routine and a linearised model at the chosen operating point.
  \item The control structure under test, with every gain named and recorded.
\end{itemize}

\subsection{Required outputs}

\begin{itemize}
  \item Time histories of position, altitude, airspeed, attitude and angular rates.
  \item Control surface deflections and throttle command, with the limits visible.
  \item Estimator output compared with the simulated truth state.
  \item Step and doublet responses for each loop, at more than one operating point.
  \item Trim results and the linear model's modes with their damping and frequency.
  \item Response to wind, gust and a disturbance large enough to matter.
\end{itemize}

\subsection{What the report must explain}

Do not upload plots without interpretation. State the objective, the
assumptions, the parameter sources, the solver and sample time, the expected
behaviour, the observed behaviour and the conclusion. If the model behaves
incorrectly, record the failure and the evidence used to diagnose it. A failed
simulation with a clear explanation is more useful than an unexplained
successful screenshot. Never present a tuned gain set without the model, the
operating point and the cost function or criterion it was tuned against.

\begin{callout}{SIMULATION EXIT CRITERION}
  The aircraft can be trimmed, linearised and flown in simulation; each control
  loop has a documented step or doublet response; the modes are identified and
  a poorly damped mode is explained; and a wind or disturbance case has been
  run and interpreted. A second member can rerun the model from the committed
  files and obtain the same result.
\end{callout}

\section{Reference stack: PX4 and ArduPilot}

A mature open-source autopilot stack provides the estimator, the control
cascade, the flight-mode state machine, the failsafe logic, the parameter
system, the logging system, the ground-station link and the simulation
harness. It allows the team to fly and measure a real aircraft without first
writing an entire autopilot, and it gives a reference to compare against when
in-house functions are developed.

\begin{callout}{IMPORTANT DISTINCTION}
  Using a reference stack does not remove the need to understand the control
  system. A member must be able to explain what the stack's loops do, why a
  parameter has the value it has, and what the log shows, before tuning it. A
  stack is a baseline to be understood and validated, not a black box that
  makes the aircraft fly.
\end{callout}

\subsection{Choose one baseline and record why}

PX4 and ArduPilot are both capable fixed-wing stacks with different
architectures, parameter systems and tuning workflows. Choose one as the
team's flight baseline and state the reason. Do not run two stacks on the same
airframe as an experiment during a flight campaign, and do not switch stacks
mid-campaign without retesting everything from the bench upwards.

\subsection{Record the exact version and configuration first}

Before following any tutorial, record the stack and version, the airframe
configuration selected, the firmware build, the ground-station version, the
parameter file, and the sensor and actuator mapping. Parameter names and
default values differ between releases and between vehicle types. Treat a video
or a blog post as a conceptual demonstration, and use documentation that
matches the exact version for actual parameter changes.

\subsection{Recommended bring-up workflow}

\begin{enumerate}
  \item Bring up software-in-the-loop simulation. Fly the simulated aircraft, change the wind, and confirm the log records what you expect.
  \item Move to hardware-in-the-loop: the real flight controller running the real firmware, with simulated sensors and actuators. Confirm timing, mode changes and failsafe behaviour.
  \item Power the aircraft on the bench with the propeller removed. Verify servo direction, travel, neutral, trim and mode switching, and confirm every failsafe triggers the intended action.
  \item Check the estimator on the bench: attitude follows the aircraft, heading is plausible, airspeed reads correctly with a controlled pressure source, and health flags behave.
  \item Configure conservative limits for speed, height, bank angle, pitch angle, climb rate, throttle and geofence before the first powered test.
  \item Run the first powered ground test at low energy with the aircraft restrained, following the team's safety procedure.
  \item Plan the first flight as a short, manual, low-risk flight whose purpose is to gather data, not to demonstrate autonomy.
  \item Log everything: estimator state, controller targets and outputs, actuator commands, mode transitions, battery, link quality and failsafe events.
\end{enumerate}

\subsection{What members should learn from the reference stack}

\begin{itemize}
  \item How a real estimator fuses sensors, reports health and degrades.
  \item How the control cascade is split into rate, attitude, energy and path loops, and which parameters shape each one.
  \item How flight modes and failsafe behaviour are structured as a state machine.
  \item How parameters are defined, validated and stored, and why a parameter set is part of the deliverable.
  \item How logs are recorded, what is in them and how a flight is reconstructed afterwards.
  \item What the stack does not do for the team: mission-specific requirements, competition rules, in-house research functions and their validation.
\end{itemize}

\subsection{Required reference-stack evidence}

\begin{itemize}
  \item Stack, firmware version, airframe configuration, ground-station version and the complete parameter file.
  \item A reproducible simulation or HIL case with its log and a short explanation of the observed behaviour.
  \item Bench evidence that actuator directions, travel, mode switching and every failsafe behave as intended.
  \item A tuning record stating the parameter changed, the previous and new value, the test condition and the observed effect.
  \item A conclusion explaining what the stack configuration proves and what it does not prove about the aircraft's mission performance.
\end{itemize}

\begin{warningbox}{REFERENCE-STACK EXIT CRITERION}
  The team has one documented baseline stack and parameter set; simulation and
  HIL reproduce; the bench checks pass with the propeller removed; every
  failsafe has been triggered deliberately and observed; and the log from each
  test can be reconstructed and explained by a member who was not present.
\end{warningbox}

\section{Flight-controller hardware and interfaces}

The flight controller is one board in a system, and the system is what flies.
Hardware documentation must make the whole chain visible: power, sensors,
actuators, links and logging.

\subsection{Required hardware blocks}

\subsubsection{Flight controller and timing}
Document the autopilot board and revision, the firmware target, the processor
and memory margin, the sensor set on the board, the available buses, the
logging medium, and the update rates that the control loops and estimator
actually run at.

\subsubsection{Power}
Document the avionics power architecture, the voltage rails and their limits,
redundancy or lack of it, the current budget with margin, the effect of a
propulsion brownout on the flight controller, and how the system behaves at
low battery.

\subsubsection{Sensors}
Document the inertial measurement unit, magnetometer, barometer, GNSS receiver,
airspeed sensor and any external sensors, including their mounting location,
orientation, vibration isolation, calibration and known failure modes.

\subsubsection{Actuator interfaces}
Document the servo and ESC outputs, their pulse or protocol definition, their
power source, their direction and travel limits, and the consequences of an
output loss or a stalled servo.

\subsubsection{Radio, telemetry and datalinks}
Document the RC link, the telemetry link, their ranges and failure behaviour,
the antenna placement, the failsafe configuration on both sides, and the
procedure to regain control.

\subsubsection{Ground station and operator interface}
Document the ground-station software and version, the parameter and mission
upload procedure, the pre-flight checklist presented to the operator, and how a
mode change or a failsafe is announced to the pilot.

\subsubsection{Logging and telemetry bandwidth}
Document what is logged at what rate, the storage and download procedure, and
what is telemetered live. Data that is not logged cannot be analysed after a
flight, and a link that is saturated is a safety problem.

\subsection{Component and configuration record}

For each significant item, record the selected part, the requirement it must
meet, the margin, the relevant datasheet pages, the alternatives considered,
the reason for the choice, the risks and the validation plan. Absolute maximum
ratings are not operating targets; design against recommended operating
conditions and the expected environment, including vibration, heat and
electrical noise.

\section{Software architecture, flight modes and failsafe}

The software documentation must describe the timing and the relationship
between a command and the physical response. A reader should be able to trace a
pilot or mission command through mode logic, guidance, the control cascade,
allocation, actuator output, and the checks that can interrupt it.

\subsection{Minimum topics}

\begin{itemize}
  \item Task structure and timing: what runs at what rate, and what is the deadline for each loop.
  \item Estimator input validation, timestamping and health reporting.
  \item Guidance and path-following logic, including how a path is entered and left.
  \item The control cascade, its limits, its anti-windup behaviour and its saturation handling.
  \item Control allocation, including trim offsets and actuator failure handling.
  \item The flight-mode state machine: which modes exist, what each one does, and what can switch between them.
  \item Failsafe definitions: RC loss, telemetry loss, estimator failure, battery low, geofence breach, and what each one commands.
  \item Parameter definitions with units, valid ranges and defaults.
  \item Logging: message set, rates, and how a flight is reconstructed.
  \item The interface to the propulsion, power and perception subteams, including units and message timing.
\end{itemize}

\subsection{Function-level documentation}

Important functions should state their purpose, inputs, outputs, units, valid
range, call rate, side effects and safety assumptions. Constants such as
control rates, limits, look-ahead distances, time constants and unit
conversions must have one authoritative definition. Changes to these constants
must be reviewed, because they change the physical behaviour of the aircraft.

\begin{callout}{SOFTWARE EXIT CRITERION}
  Every mode has a documented purpose and entry condition; every failsafe has
  been triggered in simulation or on the bench and its action observed; the
  parameter set is version-controlled and reproducible; and a flight log can be
  traced from a pilot input to an actuator output by reading the code and the
  log together.
\end{callout}

\section{Ground and bench testing, and safety}

The first objective is deliberately narrow: verify that the aircraft's control
chain responds correctly and safely before it ever leaves the ground. Test one
assumption at a time and capture the evidence.

\subsection{Minimum safety controls}

\begin{itemize}
  \item Remove the propeller for every bench test unless an approved guarded test specifically requires it.
  \item Follow the team's flight-approval procedure, and never conduct a first flight without an approved plan and a designated pilot.
  \item Verify failsafe behaviour deliberately, on the ground, before relying on it in the air.
  \item Confirm the geofence, height limit, distance limit and return behaviour and make sure the operator can see and reach the aircraft.
  \item Check control directions and mode switching with the aircraft restrained and the propulsion disabled or limited.
  \item Use a charged, inspected battery, a current-limited or instrumented supply where practical, and confirm the avionics rail holds up under load.
  \item Do not fly with an unresolved estimator, calibration or GPS-health warning, and do not fly in wind or visibility conditions outside the approved limits.
  \item Stop immediately for unexpected control response, oscillation, loss of link, estimator degradation or any unexplained change in behaviour.
\end{itemize}

\subsection{Bench test sequence}

\begin{enumerate}
  \item Power and communications: the flight controller boots, the ground station connects, sensors are detected and no pre-arm error is present.
  \item Estimator: attitude tracks the aircraft as it is moved by hand, heading is correct, airspeed reads correctly against a known pressure, and health flags appear and clear as expected.
  \item Actuators: direction, travel, neutral, trim and rate are correct for every output, with the propeller removed.
  \item Modes: every flight mode is selected on the ground and the resulting control behaviour is observed and recorded.
  \item Failsafe: the RC link and the telemetry link are interrupted deliberately and the aircraft's response is recorded.
  \item Logging: a log is recorded, downloaded and opened, and the expected messages are present at the expected rates.
\end{enumerate}

\subsection{Test record}

\begin{itemize}
  \item Test ID, date, operator, safety observer and objective.
  \item Airframe revision, firmware commit, parameter file, stack version and instrument list.
  \item Configuration, battery state, payload, propeller status and environmental conditions.
  \item Procedure, expected result, measured result, plots or log excerpts, and a pass or fail decision.
  \item Unexpected behaviour, likely cause, follow-up action and a linked issue or design note.
\end{itemize}

\begin{warningbox}{FIRST POWERED MILESTONE}
  With the aircraft restrained, the control chain responds correctly in every
  mode, all control directions are verified, every failsafe has been triggered
  and observed, the log contains the expected data, and the result is
  reproducible by another member.
\end{warningbox}

\section{Flight testing and flight-data review}

Flight testing is how the model is confronted with reality. The purpose of the
first flights is to collect data safely, not to demonstrate the mission.

\subsection{Planning a flight test}

\begin{itemize}
  \item State the objective in one sentence, and the single change or question the flight is meant to answer.
  \item Define the flight card: modes to be used, the sequence, the abort criteria and the recovery procedure.
  \item Confirm the aircraft configuration, battery, payload, centre of gravity and calibration state.
  \item Confirm the site, the airspace and radio requirements, the weather limits and the number of people required.
  \item Define the envelope for the day and expand it only after the data from the previous flight has been reviewed.
  \item Agree the roles of pilot, observer and data operator before the flight, not during it.
\end{itemize}

\subsection{Envelope expansion}

Expand the flight envelope in small, documented steps: one change at a time,
reviewed before the next flight. Do not combine a new mode, a new gain set and
a new mission in a single flight. If a flight produces an unexpected result,
stop expanding the envelope until the cause is understood.

\subsection{Flight-data review}

Review the log against the flight card before the next flight. A useful review
answers: did the aircraft do what was commanded; did any limit or failsafe
activate; did the estimator degrade; was the control response consistent with
the simulation; and what is the single most important change to make next. Record
the review as a short document, not only as a conversation.

\subsection{Required flight evidence}

\begin{itemize}
  \item Flight card, objective, configuration, firmware and parameter version.
  \item The complete raw log for the flight, plus a processed summary with the axes labelled and units stated.
  \item Commanded against measured attitude, airspeed, height and path, with the error visible.
  \item Actuator commands and saturation events.
  \item Estimator health, innovation and any reset during the flight.
  \item Failsafe, mode transition and limit events with timestamps.
  \item A conclusion separating what the flight proves, what it does not prove, and what is still uncertain.
\end{itemize}

\begin{callout}{FLIGHT-TEST EXIT CRITERION}
  The aircraft completes the planned flight card; commanded and measured values
  agree within the stated tolerance; no unexplained event appears in the log;
  the data is archived with the configuration that produced it; and the next
  flight's objective follows from this flight's data.
\end{callout}

\Resource
  {ArduPilot: first flight and landing}
  {https://ardupilot.org/plane/docs/first-flight-landing-page.html}
  {A structured first-flight and landing procedure, including pre-flight checks and what to observe in the log.}

\Resource
  {PX4: flight log analysis}
  {https://docs.px4.io/main/en/log/flight_log_analysis.html}
  {How flight logs are recorded, downloaded and reviewed, including the tools used for plotting and comparison.}

\Resource
  {PlotJuggler}
  {https://plotjuggler.io/}
  {Fast time-series plotting for flight logs; useful for comparing commanded and measured signals across a whole flight.}

\section{Guidance, mission autonomy and team interfaces}

Guidance turns a mission into a path and a set of references the control loops
can follow. It should be developed only after the inner loops are validated,
because a guidance law cannot compensate for an aircraft that does not hold
attitude or airspeed.

\subsection{Concepts to learn}

\begin{itemize}
  \item Waypoint navigation, line and arc following, and how a path is entered and exited.
  \item Look-ahead based path following, and the effect of its distance parameter at different speeds.
  \item Total energy management of speed and height, and how a commanded climb changes the speed reference.
  \item Mission modes: take-off, cruise, loiter, return, approach and landing, and what each one must guarantee.
  \item Flight-data review of a mission: path tracking error, energy behaviour and mode transitions.
  \item The boundary between conventional guidance and learning-based or optimisation-based autonomy, and why the interface between them must be defined before either is developed.
\end{itemize}

\subsection{Required guidance evidence}

\begin{itemize}
  \item A documented mission definition with waypoints, altitudes, speeds and the expected behaviour at each leg.
  \item Simulation results showing path tracking error for at least one disturbed case.
  \item The commanded and achieved values for speed, height and bank angle, with the limits shown.
  \item A statement of what the guidance layer does when the mission cannot be completed safely.
\end{itemize}

\subsection{Interfaces to the other subteams}

\begin{itemize}
  \item \textbf{Motor control and ESC:} the throttle or thrust command, its units, its rate limit, its failsafe value, and the telemetry the controller may consume.
  \item \textbf{Battery management and power:} pack state, available current, voltage limits and the low-battery policy the controller must enforce.
  \item \textbf{Sensors and estimation:} which states are provided, at what rate, with what health indication, and what the controller does when a state is lost.
  \item \textbf{Vision and autonomy:} what the perception layer outputs, in which frame, at what rate, with what confidence, and the rule that a perception failure must degrade to a safe conventional behaviour rather than take over flight control.
  \item \textbf{Structures and mechanical:} mass, centre of gravity, inertia, control-surface limits and actuator installation, since all of these appear directly in the model.
  \item \textbf{Ground station and operations:} the pre-flight checklist, the flight card and the data archive, because a control result is only usable if the operation that produced it is repeatable.
\end{itemize}

\begin{callout}{AUTONOMY ENTRY CRITERION}
  The inner loops are validated in flight, the estimator is trustworthy, every
  failsafe has been demonstrated, and the interface to whatever provides the
  higher-level decision is defined in units, frames, rates and failure
  behaviour. Autonomy must not be used to hide unverified low-level behaviour.
\end{callout}

\section{Repository structure}

The repository should make the current design state easy to find. The main
README is an entry point; detailed technical work belongs in focused Markdown
documents and version-controlled project files.

\begin{tcolorbox}[
  breakable,
  colback=DARESPaper,
  colframe=DARESLine,
  boxrule=0.4pt,
  arc=0pt,
  left=3mm,right=3mm,top=3mm,bottom=3mm]
\RepoEntry{README.md}{project overview, status, and quick start}
\RepoEntry{docs/flight-mechanics/}{learning notes and reviewed explanations}
\RepoEntry{docs/control-theory/}{loop structure, tuning notes, and analysis}
\RepoEntry{docs/design-notes/}{engineering decisions and trade-offs}
\RepoEntry{docs/interfaces/}{interface definitions shared with other subteams}
\RepoEntry{docs/datasheets/}{links or approved copies of key datasheets}
\RepoEntry{model/}{airframe model, aerodynamic data, trim and linearisation}
\RepoEntry{model/aero-data/}{coefficient sources, wind-tunnel or CFD results, and assumptions}
\RepoEntry{simulation/}{SITL and HIL cases, parameters, scripts, and reports}
\RepoEntry{autopilot/params/}{versioned parameter sets and the reason for each change}
\RepoEntry{autopilot/scripts/}{analysis, log processing, and automation}
\RepoEntry{hardware/}{airframe interface, wiring, sensor mounting, and review notes}
\RepoEntry{testing/bench/}{bench procedures, results, and logs}
\RepoEntry{testing/flight/}{flight cards, logs, reviews, and conclusions}
\RepoEntry{results/}{approved summary plots and milestone evidence}
\end{tcolorbox}

\subsection{File and commit expectations}

\begin{itemize}
  \item Use descriptive filenames and include units or dates where they prevent ambiguity.
  \item Do not upload an unexplained screenshot, log, model or binary file.
  \item Link research sources and record the exact document version or datasheet revision used.
  \item Record failures as well as successes; do not remove evidence because a test failed.
  \item Keep generated exports separate from editable source files, and keep raw logs separate from processed results.
  \item Version parameter sets with the flight or test that produced them; never overwrite the set that flew.
  \item Use issues or design notes for unresolved questions and link the relevant commit, log or test report.
  \item Request review for changes affecting control-surface direction, limits, failsafe behaviour, estimator configuration or anything that changes the aircraft's physical response.
\end{itemize}

\section{Required document templates}

\subsection{Simulation note}
\begin{itemize}
  \item Objective and question being tested.
  \item Model scope, assumptions, parameter sources, trim point, solver and sample time.
  \item Inputs and expected behaviour.
  \item Results with units and explanation.
  \item Limitations, problems, conclusion and next action.
\end{itemize}

\subsection{Design decision record}
\begin{itemize}
  \item Decision and date.
  \item Requirement or problem being addressed.
  \item Alternatives considered and evidence reviewed.
  \item Selected option, calculations, margins, benefits and disadvantages.
  \item Risks, validation plan, owner and review status.
\end{itemize}

\subsection{Bench test report}
\begin{itemize}
  \item Test ID, objective, operator, safety observer, date and location.
  \item Airframe and wiring revision, firmware commit, parameter file and instruments.
  \item Safety controls and pre-test checks.
  \item Procedure, expected result, measured data and uncertainty where relevant.
  \item Pass or fail decision, fault observations, conclusion and linked follow-up issue.
\end{itemize}

\subsection{Flight card and flight report}
\begin{itemize}
  \item Flight ID, objective, site, date, pilot, observer and data operator.
  \item Configuration, battery state, centre of gravity, firmware and parameter version.
  \item Modes and sequence, abort criteria and recovery procedure.
  \item Result against the objective, log reference, events and anomalies.
  \item Conclusion, follow-up actions and the objective of the next flight.
\end{itemize}

\subsection{Datasheet review}

Before using a component, review at least the absolute maximum ratings,
recommended operating conditions, pin descriptions, electrical
characteristics, timing requirements, thermal information and typical
application circuit. Record which limits are design requirements and which must
be validated during testing.

\section{First tasks for new members}

These tasks turn passive learning into reviewable engineering output. New
members should work with an experienced member, but they should still form and
test their own explanation.

\subsection{Task 1 - Principle note}
Write a short explanation of how lift, drag, moments and the four control
channels interact to make a fixed-wing aircraft follow a commanded path.
Include the effect of a change in airspeed and of a bank angle.

\subsection{Task 2 - Model and trim}
Take the team's airframe data, trim the model at a stated speed and height,
linearise it, and report the modes with their frequency and damping. Say which
coefficient you trust least and why.

\subsection{Task 3 - Simulation}
Reproduce a simulation case: a step or doublet in one channel, plus one
disturbed case with wind. Explain the response, the limits that were reached,
and what you would change in the control structure.

\subsection{Task 4 - Reference-stack bring-up}
Record the exact stack and version, run the aircraft in SITL, complete the
bench tests with the propeller removed, and document every failsafe you
triggered and what the aircraft did.

\subsection{Task 5 - Interface review}
Choose one interface from Section 9 - propulsion, power, estimation,
perception, structures or operations - and document its units, rates, limits
and failure behaviour, with one recommendation.

\subsection{Task 6 - GitHub submission and short presentation}
Upload the editable source, the readable result, the parameter or configuration
file, and a Markdown explanation. Link every source, state what remains
uncertain, and present one mistake or unexpected result and how you investigated
it.

\begin{callout}{REVIEW STANDARD}
  The goal is not to appear correct on the first attempt. The goal is to make
  assumptions visible, test them, record the evidence, and improve the next
  design. A control result is only as good as the configuration and data that
  produced it.
\end{callout}

\section{Milestones and definition of done}

\subsection{Milestone 1 - Fundamentals}
A member can explain the forces, moments, control channels and reference frames
of a fixed-wing aircraft without copying a diagram, and can say what each
control surface does and why.

\subsection{Milestone 2 - Model, trim and linearisation}
The six-degree-of-freedom model runs, its parameters and sources are
documented, a trim solution exists, the linear model's modes are identified,
and the results are explained.

\subsection{Milestone 3 - Simulation and control structure}
The control cascade runs in simulation; each loop has a documented step or
doublet response; one disturbed case is analysed; and gains are recorded with
the conditions that produced them.

\subsection{Milestone 4 - Reference stack in SITL and HIL}
The baseline stack and parameter set are documented; SITL reproduces; HIL
confirms timing, modes and failsafe behaviour; and logs are archived.

\subsection{Milestone 5 - Bench validation}
With the propeller removed, every actuator direction and limit is verified,
every mode is selected, every failsafe is triggered and observed, and the
aircraft is declared ready for a first flight by the responsible member.

\subsection{Milestone 6 - First flight and envelope expansion}
The aircraft flies the planned card; commanded and measured values agree within
the stated tolerance; no unexplained event appears in the log; and the envelope
is expanded only after each flight is reviewed.

\subsection{Milestone 7 - Mission modes and guidance}
Waypoint navigation, loiter, return and approach behaviour are validated in
simulation and then in flight, with path tracking error and energy behaviour
documented.

\subsection{Milestone 8 - Autonomy readiness}
The inner loops are trustworthy, the estimator is validated, failsafe behaviour
is demonstrated, the interface to the higher-level decision layer is defined,
and the team has an agreed validation plan for autonomous functions.

\subsection{Project-level definition of done}
\begin{itemize}
  \item Editable model, configuration and parameter sources are committed with build or run instructions.
  \item Requirements, assumptions, aerodynamic data sources and safety limits are documented.
  \item Simulation, bench and flight evidence is stored with units, versions and the configuration that produced it.
  \item Known limitations and unresolved risks are explicit.
  \item Another team member can reproduce the result without relying on the original author.
  \item A reviewer has approved the milestone against its stated exit criteria.
\end{itemize}

\section{Learning resources}

Use these resources as starting points, then confirm implementation details
against the relevant airframe, autopilot, sensor and actuator documentation.
Video explanations are useful for intuition but are not the final authority for
flight limits or safety decisions.

\Resource
  {PX4 user guide: fixed-wing configuration}
  {https://docs.px4.io/main/en/config_fw/}
  {How a fixed-wing vehicle is configured in PX4, from airframe selection and control-surface mapping to tuning.}

\Resource
  {PX4 user guide: fixed-wing flight modes}
  {https://docs.px4.io/main/en/flight_modes_fw/}
  {Definitions of manual, stabilised, altitude, position, mission and return modes, and what each one commands.}

\Resource
  {PX4 user guide: simulation}
  {https://docs.px4.io/main/en/simulation/}
  {Software-in-the-loop and hardware-in-the-loop simulation setup, including jMAVSim, Gazebo and FlightGear back ends.}

\Resource
  {ArduPilot developer documentation}
  {https://ardupilot.org/dev/}
  {Architecture, build system and code structure of a mature open-source autopilot; use it to read how the control loops are actually implemented.}

\Resource
  {ArduPilot Plane documentation}
  {https://ardupilot.org/plane/}
  {The fixed-wing side of ArduPilot: setup, tuning, modes, failsafe and parameter reference.}

\Resource
  {Small Unmanned Aircraft: Theory and Practice - companion code}
  {https://github.com/randybeard/uavbook}
  {Simulation and control code accompanying the standard fixed-wing UAV textbook; a good source of implementable guidance and autopilot structure.}

\Resource
  {BYU MAGICC Lab UAV book site}
  {https://uavbook.byu.edu/}
  {Course material and chapter resources matching the fixed-wing UAV textbook.}

\Resource
  {MIT OpenCourseWare 16.30 Feedback Control Systems}
  {https://ocw.mit.edu/courses/16-30-feedback-control-systems-fall-2010/}
  {State-space control foundations: modelling, linearisation, observability, state feedback and observers.}

\Resource
  {Brian Douglas - control system lectures}
  {https://engineeringmedia.com/}
  {Intuitive explanations of PID, loop shaping, state-space control and Kalman filtering.}

\Resource
  {MathWorks Aerospace Blockset}
  {https://www.mathworks.com/help/aeroblks/}
  {Reference blocks for six-degree-of-freedom motion, atmosphere, aerodynamics and flight control, with documented conventions.}

\Resource
  {MathWorks UAV Toolbox}
  {https://www.mathworks.com/help/uav/}
  {Tools for UAV scenario simulation, coordinate transformations and flight-log analysis.}

\Resource
  {JSBSim flight dynamics model}
  {https://jsbsim.sourceforge.net/}
  {Open-source six-degree-of-freedom flight dynamics engine used by several simulators; useful for cross-checking an in-house model.}

\Resource
  {ArduPilot Methodic Configurator}
  {https://github.com/ArduPilot/MethodicConfigurator}
  {A structured, documented workflow for configuring a new ArduPilot aircraft parameter by parameter, with reasons recorded.}

\Resource
  {MAVLink developer guide}
  {https://mavlink.io/en/}
  {The message protocol between autopilot, ground station and companion computers; read it before defining an interface.}

\Resource
  {PX4 ULog file format}
  {https://docs.px4.io/main/en/dev_log/ulog_file_format.html}
  {The structure of a PX4 flight log, needed when writing log-processing scripts.}

\Resource
  {pyulog - PX4 log parsing in Python}
  {https://github.com/PX4/pyulog}
  {Python library for reading ULog files into dataframes for analysis and plotting.}

\subsection{Useful search terms}
\begin{itemize}
  \item fixed-wing flight dynamics stability and control
  \item aircraft equations of motion six degree of freedom
  \item trim and linearisation longitudinal lateral modes
  \item short period phugoid Dutch roll spiral mode
  \item total energy control system TECS fixed wing
  \item L1 adaptive path following UAV
  \item control allocation aircraft
  \item extended Kalman filter attitude estimation UAV
  \item pitot static airspeed calibration UAV
  \item PX4 fixed-wing tuning guide
  \item ArduPilot plane tuning and failsafe
  \item software in the loop and hardware in the loop autopilot
\end{itemize}

\subsection{Start this week}
\begin{itemize}
  \item Read the flight-mechanics and reference-frame sections of this guide and write a one-page explanation in your own words.
  \item Open the team's simulation and run one case; change a single parameter and explain what changed.
  \item Bring up the reference stack in SITL and fly one simulated flight, then find that flight in the log.
  \item Review one interface in Section 9 and write down what is undefined about it.
  \item Upload your editable source and evidence, then present one unresolved question and the next test that would reduce the uncertainty.
\end{itemize}

\end{document}
